
Across N=296 open-weight LLMs from the Open LLM Leaderboard v1, MMLU explains 49.3\% of TruthfulQA MC2 variance and 76.3\% of bias-proxy (ARC Challenge) variance. Controlling for MMLU via partial Spearman reduces the pairwise alignment benchmark correlation from 0.732 to 0.343 --- a 53.1\% reduction --- with Fisher z = 6.97, p = 3.22 $\times 10^{-12}$, and non-overlapping BCa confidence intervals. The reduction confirms that scale confounding accounts for more than half of the raw positive correlation between these two benchmarks. A positive residual partial correlation (0.343, p = 1.40 $\times 10^{-9})$ indicates non-zero alignment-specific co-movement beyond scale, though the structural classification of this residual remains ambiguous at N=296, as the BCa CI [0.180, 0.492] spans the pre-specified coherence threshold of 0.40.

The bias benchmark used throughout is ARC Challenge accuracy as a proxy for HELM Lite BBQ per-model scores, which were unavailable in the execution environment. All secondary claims regarding factuality-bias coupling are qualified by this substitution. The safety dimension (HarmBench) could not be evaluated due to zero model overlap between datasets. An optional RLHF sign test (N=300 base/chat pairs) yields a null result (p = 0.686), providing no cross-model evidence that RLHF fine-tuning systematically improves bias-proxy scores.

The most immediate next steps are replication with genuine HELM Lite BBQ per-model accuracy scores and extension to a larger model cohort (LLM Leaderboard v2, N > 1,000) to resolve the scenario classification ambiguity. The Fisher z raw-vs-partial difference test, as demonstrated here, represents a low-cost diagnostic that should become standard practice before interpreting alignment benchmark co-movement as evidence of aligned training objectives.

